% from ega3-7-fr.tex
\oldpage[III]{201}
à partir de $T_\bullet^A$ par extension des scalaires à $A'$, et il est clair
alors que a$'$) entraîne que (7.7.5.3) est un isomorphisme. Inversement, supposons
toujours $Y=\operatorname{Spec}(A)$ affine et soit $A'$ la $A$-algèbre
$A\oplus M$, où $M$ est un $A$-module quelconque, la multiplication dans $A'$
étant donnée par
\[
  (a_1,m_1)(a_2,m_2)
  =
  (a_1a_2,a_1m_2+a_2m_1);
\]
alors
\[
  T_p^{A'}(A')
  =
  T_p(A\oplus M)
  =
  T_p(A)\oplus T_p(M),
\]
et l'hypothèse que (7.7.5.3) soit bijectif entraîne qu'il en est de même de
l'application canonique
\[
  T_p(A)\otimes_A M\longrightarrow T_p(M),
\]
autrement dit d) entraîne a$''$), ce qui achève la démonstration.
\end{proof}

\begin{theorem}[7.7.6]
Soient $Y$ un préschéma localement noethérien, $f:X\to Y$ un morphisme propre,
$\mathscr F$ un $\mathscr O_X$-Module cohérent et $Y$-plat. Il existe alors un
$\mathscr O_Y$-Module cohérent $\mathscr Q$ (déterminé à isomorphisme unique près)
et un isomorphisme de foncteurs en le $\mathscr O_Y$-Module quasi-cohérent
$\mathscr M$ :
\[
  f_*(\mathscr F\otimes_{\mathscr O_Y}\mathscr M)
  \xrightarrow{\sim}
  \mathscr{H}\!om_{\mathscr O_Y}(\mathscr Q,\mathscr M)
  \tag{7.7.6.1}\label{III.7.7.6.1.fr}
\]
(d'où un isomorphisme de foncteurs
\[
  \Gamma(X,\mathscr F\otimes_{\mathscr O_Y}\mathscr M)
  \xrightarrow{\sim}
  \operatorname{Hom}_{\mathscr O_Y}(\mathscr Q,\mathscr M)).
  \tag{7.7.6.2}\label{III.7.7.6.2.fr}
\]
\end{theorem}

\begin{proof}
En effet, comme $\mathscr M\mapsto\mathscr F\otimes_{\mathscr O_Y}\mathscr M$ est
exact (0\textsubscript{I}, 6.7.4) et $f_*$ exact à gauche, le foncteur
$\mathscr M\mapsto f_*(\mathscr F\otimes_{\mathscr O_Y}\mathscr M)$ est exact à
gauche. Il suffit alors d'appliquer l'équivalence de (7.7.5, b)) et (7.7.5,
b$'$)) pour $p=1$.
\end{proof}

\begin{corollary}[7.7.7]
Soient $Y$ un préschéma localement noethérien, $f:X\to Y$ un morphisme propre,
$\mathscr F$, $\mathscr F'$ deux $\mathscr O_X$-Modules cohérents et $Y$-plats,
$u:\mathscr F\to\mathscr F'$ un homomorphisme. Considérons les deux foncteurs en
le $\mathscr O_Y$-Module quasi-cohérent $\mathscr M$ :
\begin{align*}
  \mathscr T(\mathscr M)
  &=
  \operatorname{Ker}
  \bigl(
    f_*(\mathscr F\otimes_{\mathscr O_Y}\mathscr M)
    \longrightarrow
    f_*(\mathscr F'\otimes_{\mathscr O_Y}\mathscr M)
  \bigr),\\
  T(\mathscr M)
  &=
  \Gamma(Y,\mathscr T(\mathscr M))\\
  &=
  \operatorname{Ker}
  \bigl(
    \Gamma(X,\mathscr F\otimes_{\mathscr O_Y}\mathscr M)
    \longrightarrow
    \Gamma(X,\mathscr F'\otimes_{\mathscr O_Y}\mathscr M)
  \bigr).
\end{align*}
Alors il existe un $\mathscr O_Y$-Module cohérent $\mathscr R$ (déterminé à
isomorphisme unique près) et des isomorphismes de foncteurs
\[
  \mathscr T(\mathscr M)
  \xrightarrow{\sim}
  \mathscr{H}\!om_{\mathscr O_Y}(\mathscr R,\mathscr M)
  \tag{7.7.7.1}\label{III.7.7.7.1.fr}
\]
\[
  T(\mathscr M)
  \xrightarrow{\sim}
  \operatorname{Hom}_{\mathscr O_Y}(\mathscr R,\mathscr M).
  \tag{7.7.7.2}\label{III.7.7.7.2.fr}
\]
\end{corollary}

\begin{proof}
On peut se borner à démontrer (7.7.7.2) ; cela prouvera en effet (7.7.7.1)
dans le cas où $Y$ est affine, et on passera de là au cas général en raisonnant
comme dans la démonstration de l'équivalence de (7.7.5, b) et b$'$)), grâce à
l'unicité à isomorphisme unique près d'un représentant d'un foncteur
représentable (0, 8.1.8). Il résulte de (7.7.6) qu'il existe deux
$\mathscr O_Y$-Modules cohérents $\mathscr Q$, $\mathscr Q'$ définissant des
isomorphismes fonctoriels
\[
  \Gamma(X,\mathscr F\otimes_{\mathscr O_Y}\mathscr M)
  \xrightarrow{\sim}
  \operatorname{Hom}_{\mathscr O_Y}(\mathscr Q,\mathscr M),
  \qquad
  \Gamma(X,\mathscr F'\otimes_{\mathscr O_Y}\mathscr M)
  \xrightarrow{\sim}
  \operatorname{Hom}_{\mathscr O_Y}(\mathscr Q',\mathscr M).
\]
Or, $u:\mathscr F\to\mathscr F'$ définit canoniquement un morphisme de foncteurs
\[
  \Gamma(X,\mathscr F\otimes_{\mathscr O_Y}\mathscr M)
  \longrightarrow
  \Gamma(X,\mathscr F'\otimes_{\mathscr O_Y}\mathscr M);
\]

